\begin{abstract}
A coding agent edits files and executes shell commands with its developer’s privileges, allowing malicious requests to translate directly into harmful actions or functional malware. Existing defenses have complementary limitations: weight-level alignment is unavailable to API-only deployers, whereas input filters and execution-boundary monitors require auxiliary classification or checking components along the agent’s trajectory. We therefore introduce \sysname, a system-prompt defense that synthesizes \emph{security skills} offline from known attacks or recorded agent failures. These skills are injected into the system prompt at session start and remain active throughout the tool-use loop. Unlike a reference monitor, they protect the system by defining the security policies the model should follow during execution. Due to the limited system-prompt space, we examine three fixed-budget provisioning scopes: \emph{all-classes}, with one skill covering all threat classes, \emph{per-bundle}, with one skill targeting a related subset, and \emph{per-class}, with one skill dedicated to a single known class and used as the upper-bound reference. None requires runtime request classification or routing. Across six large language models on RedCode, the default \emph{all-classes} skill reduces malware-generation severity from 3.37 to 0.58 and achieves a 43.6\% execution attack success rate, comparable to Llama Guard~3’s 42.7\% without its separate 8B classifier. The \emph{per-bundle} and class-fixed \emph{per-class} settings further reduce this rate to 36.2\% and 14.5\%, respectively. Under two non-adaptive jailbreak families, \sysname continues to outperform all baselines on malware generation. Across 731 benign task descriptions, \sysname yields a mean safety-refusal rate of 0.14\%. These results demonstrate the potential of prompt-space security skills to prevent harmful actions and malware generation for LLM coding agents.
\end{abstract}
